% GENERATED FILE. Edit canonical lessons, then run npm run generate:books.
% canonical-source-hash: fnv1a64:640f1b77e93a863a
% canonical-lessons: GU-C117-kadhi, GU-C117-bhajiya, GU-C117-catni, GU-C117-halvo, GU-C117-khir

\chapter{Kadhi, Fritters, Chutney, Halwa, Rice pudding}
\label{ch:gu-kadhi-fritters-chutney-halwa-rice-pudding}
\hlchaptermodality{117}

\begin{chapteropening}
\textbf{By the end of this chapter:} \emph{I can say kadhi, fritters, chutney, halwa and rice pudding.}

Complete the last lesson of chapter 117: I can say kadhi, fritters, chutney, halwa and rice pudding.
\end{chapteropening}

\section[\texorpdfstring{\gu{કઢી} (kaḍhī)}{કઢી (kaḍhī)}]{\gu{કઢી} (kaḍhī) — kadhi}

\label{lesson:GU-C117-kadhi}



\begin{quote}
Before the new one: say the Gujarati for peanuts, then the Gujarati for dhokla.
\end{quote}

\subsection*{You'll want to know: \gu{કઢી}}
\textbf{\gu{કઢી}} (\emph{kaḍhī}) — \textquotedblleft{}kadhi\textquotedblright{}.

\begin{grammarlens}[title={What you've built}]
One more word for this chapter.
\end{grammarlens}

\subsection*{Your turn}
\begin{itemize}
\raggedright
  \item \emph{Say it:} \emph{kaḍhī}
  \item \emph{Say it:} \emph{kaḍhī}, once more
  \item \emph{Say it:} say \emph{kaḍhī}
  \item \emph{Recall:} say the Gujarati for peanuts, then the Gujarati for dhokla, then say \emph{kaḍhī} again
\end{itemize}

\begin{tcolorbox}[breakable,colback=teal!4,colframe=teal!35!black,title={Before you move on}]
What does \gu{કઢી} mean? (\textquotedblleft{}Kadhi\textquotedblright{}.) Say it once more.
\end{tcolorbox}
\section[\texorpdfstring{\gu{ભજિયાં} (bhajiyā̃)}{ભજિયાં (bhajiyā̃)}]{\gu{ભજિયાં} (bhajiyā̃) — fritters}

\label{lesson:GU-C117-bhajiya}



\begin{quote}
Before the new one: say the Gujarati for dhokla, then the Gujarati for kadhi.
\end{quote}

\subsection*{You'll want to know: \gu{ભજિયાં}}
\textbf{\gu{ભજિયાં}} (\emph{bhajiyā̃}) — \textquotedblleft{}fritters\textquotedblright{}.

\begin{grammarlens}[title={What you've built}]
One more word for this chapter.
\end{grammarlens}

\subsection*{Your turn}
\begin{itemize}
\raggedright
  \item \emph{Say it:} \emph{bhajiyā̃}
  \item \emph{Say it:} \emph{bhajiyā̃}, once more
  \item \emph{Say it:} say \emph{bhajiyā̃}
  \item \emph{Recall:} say the Gujarati for dhokla, then the Gujarati for kadhi, then say \emph{bhajiyā̃} again
\end{itemize}

\begin{tcolorbox}[breakable,colback=teal!4,colframe=teal!35!black,title={Before you move on}]
What does \gu{ભજિયાં} mean? (\textquotedblleft{}Fritters\textquotedblright{}.) Say it once more.
\end{tcolorbox}
\section[\texorpdfstring{\gu{ચટણી} (caṭṇī)}{ચટણી (caṭṇī)}]{\gu{ચટણી} (caṭṇī) — chutney}

\label{lesson:GU-C117-catni}



\begin{quote}
Before the new one: say the Gujarati for kadhi, then the Gujarati for fritters.
\end{quote}

\subsection*{You'll want to know: \gu{ચટણી}}
\textbf{\gu{ચટણી}} (\emph{caṭṇī}) — \textquotedblleft{}chutney\textquotedblright{}.

\begin{grammarlens}[title={What you've built}]
One more word for this chapter.
\end{grammarlens}

\subsection*{Your turn}
\begin{itemize}
\raggedright
  \item \emph{Say it:} \emph{caṭṇī}
  \item \emph{Say it:} \emph{caṭṇī}, once more
  \item \emph{Say it:} say \emph{caṭṇī}
  \item \emph{Recall:} say the Gujarati for kadhi, then the Gujarati for fritters, then say \emph{caṭṇī} again
\end{itemize}

\begin{tcolorbox}[breakable,colback=teal!4,colframe=teal!35!black,title={Before you move on}]
What does \gu{ચટણી} mean? (\textquotedblleft{}Chutney\textquotedblright{}.) Say it once more.
\end{tcolorbox}
\section[\texorpdfstring{\gu{હલવો} (halvo)}{હલવો (halvo)}]{\gu{હલવો} (halvo) — halwa}

\label{lesson:GU-C117-halvo}



\begin{quote}
Before the new one: say the Gujarati for fritters, then the Gujarati for chutney.
\end{quote}

\subsection*{You'll want to know: \gu{હલવો}}
\textbf{\gu{હલવો}} (\emph{halvo}) — \textquotedblleft{}halwa\textquotedblright{}.

\begin{grammarlens}[title={What you've built}]
One more word for this chapter.
\end{grammarlens}

\subsection*{Your turn}
\begin{itemize}
\raggedright
  \item \emph{Say it:} \emph{halvo}
  \item \emph{Say it:} \emph{halvo}, once more
  \item \emph{Say it:} say \emph{halvo}
  \item \emph{Recall:} say the Gujarati for fritters, then the Gujarati for chutney, then say \emph{halvo} again
\end{itemize}

\begin{tcolorbox}[breakable,colback=teal!4,colframe=teal!35!black,title={Before you move on}]
What does \gu{હલવો} mean? (\textquotedblleft{}Halwa\textquotedblright{}.) Say it once more.
\end{tcolorbox}
\section[\texorpdfstring{\gu{ખીર} (khīr)}{ખીર (khīr)}]{\gu{ખીર} (khīr) — rice pudding}

\label{lesson:GU-C117-khir}



\begin{quote}
Before the new one: say the Gujarati for chutney, then the Gujarati for halwa.
\end{quote}

\subsection*{You'll want to know: \gu{ખીર}}
\textbf{\gu{ખીર}} (\emph{khīr}) — \textquotedblleft{}rice pudding\textquotedblright{}.

\begin{grammarlens}[title={What you've built}]
One more word for this chapter.
\end{grammarlens}

\subsection*{Your turn}
\begin{itemize}
\raggedright
  \item \emph{Say it:} \emph{khīr}
  \item \emph{Say it:} \emph{khīr}, once more
  \item \emph{Say it:} say \emph{khīr}
  \item \emph{Recall:} say the Gujarati for chutney, then the Gujarati for halwa, then say \emph{khīr} again
\end{itemize}

\begin{tcolorbox}[breakable,colback=teal!4,colframe=teal!35!black,title={Before you move on}]
What does \gu{ખીર} mean? (\textquotedblleft{}Rice pudding\textquotedblright{}.) Say it once more.
\end{tcolorbox}
